\documentclass[10pt]{amsart}
\usepackage{geometry}                % See geometry.pdf to learn the layout options. There are lots.
\geometry{letterpaper}                   % ... or a4paper or a5paper or ... 
%\geometry{landscape}                % Activate for for rotated page geometry
%\usepackage[parfill]{parskip}    % Activate to begin paragraphs with an empty line rather than an indent
\usepackage{graphicx}
\usepackage{amssymb}
\usepackage{epstopdf}
\DeclareGraphicsRule{.tif}{png}{.png}{`convert #1 `dirname #1`/`basename #1 .tif`.png}

\title{On the surface pressure equation advection diffusion equation analogy}
\author{Github Copilot}
\date{}                                           % Activate to display a given date or no date

\begin{document}
\maketitle
%\section{}
\subsection{Sign-definiteness of diffusivity}

Let
\begin{equation}
\varphi=\sqrt{\frac{i f}{\nu\,\mathrm{Ek}}},
\qquad
\mathfrak A
=
\frac{\varphi^{-1}\tanh(\varphi H)-H}{\nu \varphi^2\mathrm{Ek}},
\end{equation}
where \(H>0\), \(\nu>0\), \(\mathrm{Ek}>0\), and \(f \neq 0\).
Introducing
\begin{equation}
\sigma=\operatorname{sgn}(f),
\qquad
\alpha=\sqrt{\frac{|f|}{2\nu\,\mathrm{Ek}}},
\qquad
s=\alpha H>0,
\end{equation}
gives
\begin{equation}
\varphi=\alpha(1+i\sigma),
\qquad
\nu \varphi^2\mathrm{Ek}=if.
\end{equation}
Using
\begin{equation}
\tanh\!\left[s(1+i\sigma)\right]
=
\frac{\sinh(2s)+i\sigma\sin(2s)}
     {\cosh(2s)+\cos(2s)},
\end{equation}
we obtain
\begin{align}
\frac{\tanh(\varphi H)}{\varphi}
&=
\frac{
\sinh(2s)+\sin(2s)
+i\sigma\left[\sin(2s)-\sinh(2s)\right]
}{
2\alpha\left[\cosh(2s)+\cos(2s)\right]
},
\\
\mathfrak A_r
\equiv \operatorname{Re}\mathfrak A
&=
\frac{\sin(2s)-\sinh(2s)}
{
2\alpha |f|
\left[\cosh(2s)+\cos(2s)\right]
}.
\end{align}
For \(s>0\),
\begin{equation}
\cosh(2s)+\cos(2s)>0
\end{equation}
and
\begin{equation}
\sinh(2s)-\sin(2s)>0,
\end{equation}
the latter following, for example, from
\begin{equation}
\sinh(2s)-\sin(2s)
=
\int_0^{2s}\left(\cosh r-\cos r\right)\,dr>0.
\end{equation}
Consequently,
\begin{equation}
\boxed{\mathfrak A_r<0}
\qquad
(H>0,\ \nu>0,\ \mathrm{Ek}>0,\ f\neq0).
\end{equation}

Thus the positive-definite diffusivity is
\begin{equation}
K\equiv-\mathfrak A_r
=
\frac{\sinh(2s)-\sin(2s)}
{
2\alpha |f|
\left[\cosh(2s)+\cos(2s)\right]
}>0.
\end{equation}
The pressure-driven transport is therefore strictly down-gradient:
\begin{equation}
\boldsymbol U_{p,\mathrm{diss}}
=
\mathfrak A_r\nabla p_s
=
-K\nabla p_s.
\end{equation}
Equivalently, the second-order part of the surface-pressure equation is
\begin{equation}
\nabla\cdot\left(\mathfrak A_r\nabla p_s\right)
=
-\nabla\cdot\left(K\nabla p_s\right),
\qquad K>0.
\end{equation}

\subsection{Peclet Number}


The surface-pressure equation may be written as
\begin{equation}
-\nabla\cdot\left(K\nabla p_s\right)
+\boldsymbol{c}\cdot\nabla p_s
=S,
\end{equation}
where
\begin{equation}
K=-\mathfrak A_r>0,
\qquad
\boldsymbol{c}
=
\left(\partial_y\mathfrak A_i,-\partial_x\mathfrak A_i\right)
=
\nabla^\perp\mathfrak A_i.
\end{equation}
A natural local P\'eclet number based on a horizontal length scale \(L\) is
\begin{equation}
\boxed{
\mathrm{Pe}_L
=
\frac{L\lvert\boldsymbol{c}\rvert}{K}
=
\frac{L\lvert\nabla\mathfrak A_i\rvert}{K}
}.
\end{equation}
Consequently,
\begin{equation}
\boxed{0\leq \mathrm{Pe}_L<\infty}.
\end{equation}
In particular, \(\mathrm{Pe}_L=0\) whenever \(\mathfrak A_i\) is spatially
uniform. This occurs, for example, when both \(H\) and \(f\) are horizontally
uniform. There is no universal finite upper bound because
\(\lvert\nabla\mathfrak A_i\rvert\) may be arbitrarily large.

To examine the coefficient dependence, define
\begin{equation}
\alpha
=
\sqrt{\frac{|f|}{2\nu\,\mathrm{Ek}}},
\qquad
s=\alpha H.
\end{equation}
The positive diffusivity is
\begin{equation}
K
=
\frac{\sinh(2s)-\sin(2s)}
{
2\alpha |f|
\left[\cosh(2s)+\cos(2s)\right]
},
\end{equation}
and the imaginary part of the transport coefficient is
\begin{equation}
\mathfrak A_i
=
\frac{\operatorname{sgn}(f)}{\alpha |f|}
\left[
s
-
\frac{\sinh(2s)+\sin(2s)}
{
2\left[\cosh(2s)+\cos(2s)\right]
}
\right].
\end{equation}

If \(\mathfrak A_i\) varies by an amount comparable to
\(\lvert\mathfrak A_i\rvert\) over the horizontal scale \(L\), then
\begin{equation}
L\lvert\nabla\mathfrak A_i\rvert
\sim
\lvert\mathfrak A_i\rvert,
\end{equation}
and hence
\begin{equation}
\mathrm{Pe}_L
\sim
\frac{\lvert\mathfrak A_i\rvert}{K}.
\end{equation}
The corresponding coefficient ratio is
\begin{equation}
\frac{\lvert\mathfrak A_i\rvert}{K}
=
\frac{
\left|
2s\left[\cosh(2s)+\cos(2s)\right]
-\sinh(2s)-\sin(2s)
\right|
}{
\sinh(2s)-\sin(2s)
}.
\end{equation}

In the shallow or strongly frictional limit,
\begin{equation}
s\ll1,
\end{equation}
Taylor expansion gives
\begin{equation}
\frac{\lvert\mathfrak A_i\rvert}{K}
=
\frac{4}{5}s^2+O(s^6).
\end{equation}
Thus,
\begin{equation}
\mathrm{Pe}_L
\sim
\frac{4}{5}s^2
\ll1.
\end{equation}

In the deep or weakly frictional limit,
\begin{equation}
s\gg1,
\end{equation}
the coefficient ratio satisfies
\begin{equation}
\frac{\lvert\mathfrak A_i\rvert}{K}
=
2s-1+O\left(e^{-2s}\right).
\end{equation}
Therefore,
\begin{equation}
\mathrm{Pe}_L
\sim
2s
=
H\sqrt{\frac{2|f|}{\nu\,\mathrm{Ek}}},
\qquad s\gg1,
\end{equation}
provided that \(\mathfrak A_i\) has order-one relative variation over \(L\).
Hence the effective P\'eclet number can become very large when the water
depth is large compared with the rotational viscous boundary-layer scale:
\begin{equation}
\boxed{
\mathrm{Pe}_L
\sim
\frac{H}{\sqrt{\nu\,\mathrm{Ek}/(2|f|)}}
\gg1
}.
\end{equation}

Large depth alone, however, does not imply large advection. If \(H\) and \(f\)
are horizontally uniform, then
\begin{equation}
\nabla\mathfrak A_i=0,
\qquad
\boldsymbol{c}=0,
\qquad
\mathrm{Pe}_L=0,
\end{equation}
regardless of the value of \(s\).


\section{Inferred surface pressure velocity and diffusivity fields: redimensionalization}

%------------------------------------------------------------
% Nondimensional coefficient
%------------------------------------------------------------

Let
\begin{equation}
\widehat q
=
\sqrt{
\frac{i\widehat f}
{\widehat\nu\,\mathrm{Ek}}
},
\qquad
\widehat{\mathfrak A}
=
\frac{
\widehat q^{-1}\tanh(\widehat q\widehat H)-\widehat H
}{
\widehat\nu\,\widehat q^{\,2}\mathrm{Ek}
}.
\end{equation}

Define
\begin{equation}
\sigma=\operatorname{sgn}(\widehat f),
\qquad
\widehat\alpha
=
\sqrt{
\frac{|\widehat f|}
{2\widehat\nu\,\mathrm{Ek}}
},
\qquad
s=\widehat\alpha\widehat H,
\qquad
D(s)=\cosh(2s)+\cos(2s).
\end{equation}

Then
\begin{equation}
\widehat q
=
\widehat\alpha(1+i\sigma),
\qquad
\widehat\nu\,\widehat q^{\,2}\mathrm{Ek}
=
i\widehat f.
\end{equation}

Writing
\begin{equation}
\widehat{\mathfrak A}
=
\widehat{\mathfrak A}_r
+i\widehat{\mathfrak A}_i,
\end{equation}
the real and imaginary parts are
\begin{subequations}
\label{eq:A-real-imag-nd}
\begin{align}
\widehat{\mathfrak A}_r
&=
\frac{
\sin(2s)-\sinh(2s)
}{
2\widehat\alpha\,|\widehat f|\,D(s)
},
\label{eq:A-real-nd}
\\[1ex]
\widehat{\mathfrak A}_i
&=
\frac{
2sD(s)-\sinh(2s)-\sin(2s)
}{
2\widehat\alpha\,\widehat f\,D(s)
}.
\label{eq:A-imag-nd}
\end{align}
\end{subequations}

Equivalently,
\begin{equation}
\boxed{
\widehat{\mathfrak A}_r
=
\frac{
\sin(2s)-\sinh(2s)
}{
2\widehat\alpha\,|\widehat f|
\left[\cosh(2s)+\cos(2s)\right]
}
}
\end{equation}
and
\begin{equation}
\boxed{
\widehat{\mathfrak A}_i
=
\frac{
2s\left[\cosh(2s)+\cos(2s)\right]
-\sinh(2s)-\sin(2s)
}{
2\widehat\alpha\,\widehat f
\left[\cosh(2s)+\cos(2s)\right]
}
}.
\end{equation}

In particular,
\begin{equation}
\widehat{\mathfrak A}_r<0
\qquad
\text{for}
\qquad
\widehat H>0,\quad
\widehat\nu>0,\quad
\mathrm{Ek}>0,\quad
\widehat f\neq0.
\end{equation}

The positive nondimensional diffusivity is therefore
\begin{equation}
\widehat K
\equiv
-\widehat{\mathfrak A}_r
=
\frac{
\sinh(2s)-\sin(2s)
}{
2\widehat\alpha\,|\widehat f|
\left[\cosh(2s)+\cos(2s)\right]
}
>0.
\end{equation}


%------------------------------------------------------------
% Redimensionalization
%------------------------------------------------------------

Use the dimensional scales
\begin{equation}
x_d=L\widehat x,
\qquad
z_d=H_0\widehat z,
\qquad
H_d=H_0\widehat H,
\end{equation}
\begin{equation}
f_d=f_0\widehat f,
\qquad
\nu_d=\nu_0\widehat\nu,
\qquad
\mathrm{Ek}
=
\frac{\nu_0}{f_0H_0^2},
\end{equation}
and
\begin{equation}
p_{s,d}
=
\rho_0 f_0 L U_0\,\widehat p_s,
\qquad
\boldsymbol{\mathcal U}_d
=
H_0U_0\,\widehat{\boldsymbol{\mathcal U}},
\end{equation}
where \(\boldsymbol{\mathcal U}_d\) is the dimensional
depth-integrated velocity transport.

If
\begin{equation}
\widehat{\boldsymbol{\mathcal U}}_p
=
\widehat{\mathfrak A}\,
\widehat\nabla\widehat p_s,
\end{equation}
then
\begin{equation}
\boldsymbol{\mathcal U}_{p,d}
=
\mathfrak A_d\,\nabla_d p_{s,d},
\qquad
\mathfrak A_d
=
\frac{H_0}{\rho_0 f_0}
\widehat{\mathfrak A}.
\end{equation}

Define 
\begin{equation}
q_d
=
\sqrt{\frac{if_d}{\nu_d}},
\end{equation}
and introduce
\begin{equation}
\sigma_d=\operatorname{sgn}(f_d),
\qquad
\alpha_d
=
\sqrt{\frac{|f_d|}{2\nu_d}},
\qquad
s_d=\alpha_dH_d.
\end{equation}
The dimensional Ekman depth is $1/\alpha_d$, so $s_d$ is the fluid depth divided by the Ekman depth.


The parameter \(s\) is invariant under the redimensionalization:
\begin{equation}
s
=
\widehat\alpha\widehat H
=
\alpha_dH_d
=
s_d.
\end{equation}

The dimensional complex coefficient is
\begin{equation}
\mathfrak A_d
=
\frac{
q_d^{-1}\tanh(q_dH_d)-H_d
}{
\rho_0\nu_d q_d^{\,2}
}
=
\frac{
q_d^{-1}\tanh(q_dH_d)-H_d
}{
i\rho_0 f_d
}.
\end{equation}

Writing
\begin{equation}
\mathfrak A_d
=
\mathfrak A_{r,d}
+i\mathfrak A_{i,d},
\end{equation}
its dimensional real and imaginary parts are
\begin{subequations}
\label{eq:A-real-imag-dimensional}
\begin{align}
\mathfrak A_{r,d}
&=
\frac{
\sin(2s_d)-\sinh(2s_d)
}{
2\rho_0\alpha_d\,|f_d|
\left[\cosh(2s_d)+\cos(2s_d)\right]
},
\label{eq:A-real-dimensional}
\\[1ex]
\mathfrak A_{i,d}
&=
\frac{
2s_d\left[\cosh(2s_d)+\cos(2s_d)\right]
-\sinh(2s_d)-\sin(2s_d)
}{
2\rho_0\alpha_d\,f_d
\left[\cosh(2s_d)+\cos(2s_d)\right]
}.
\label{eq:A-imag-dimensional}
\end{align}
\end{subequations}

Equivalently,
\begin{equation}
\boxed{
\mathfrak A_{r,d}
=
\frac{
\sin(2s_d)-\sinh(2s_d)
}{
2\rho_0\alpha_d\,|f_d|
\left[\cosh(2s_d)+\cos(2s_d)\right]
}
<0
}
\end{equation}
and
\begin{equation}
\boxed{
\mathfrak A_{i,d}
=
\frac{
2s_d\left[\cosh(2s_d)+\cos(2s_d)\right]
-\sinh(2s_d)-\sin(2s_d)
}{
2\rho_0\alpha_d\,f_d
\left[\cosh(2s_d)+\cos(2s_d)\right]
}
}.
\end{equation}

The positive dimensional diffusivity coefficient is
\begin{equation}
K_d
\equiv
-\mathfrak A_{r,d}
=
\frac{
\sinh(2s_d)-\sin(2s_d)
}{
2\rho_0\alpha_d\,|f_d|
\left[\cosh(2s_d)+\cos(2s_d)\right]
}
>0.
\end{equation}

Thus the pressure-driven dimensional transport is
\begin{equation}
\boldsymbol{\mathcal U}_{p,d}
=
-K_d\nabla_d p_{s,d}
+
\mathfrak A_{i,d}\,
\boldsymbol{k}\times\nabla_d p_{s,d}.
\end{equation}

\subsection{Advecting velocity streamfunction decomposition}

The dimensional advecting streamfunction is the imaginary part of
\(\mathfrak A_d\). It separates naturally into inviscid and frictional
parts:
\begin{equation}
\mathfrak A_{i,d}
=
\mathfrak A_{i,d}^{(0)}
+
\mathfrak A_{i,d}^{(\mathrm{fr})},
\end{equation}
where
\begin{equation}
\boxed{
\mathfrak A_{i,d}^{(0)}
=
\frac{H_d}{\rho_0 f_d}
}
\end{equation}
is the classical inviscid result, and
\begin{equation}
\boxed{
\mathfrak A_{i,d}^{(\mathrm{fr})}
=
-
\frac{
\sinh(2s_d)+\sin(2s_d)
}{
2\rho_0\alpha_d f_d
\left[\cosh(2s_d)+\cos(2s_d)\right]
}
}
\end{equation}
is the frictional correction, with
\begin{equation}
\alpha_d
=
\sqrt{\frac{|f_d|}{2\nu_0}},
\qquad
s_d=\alpha_dH_d.
\end{equation}
Thus,
\begin{equation}
\boxed{
\mathfrak A_{i,d}
=
\frac{H_d}{\rho_0 f_d}
-
\frac{
\sinh(2s_d)+\sin(2s_d)
}{
2\rho_0\alpha_d f_d
\left[\cosh(2s_d)+\cos(2s_d)\right]
}
}.
\end{equation}

The corresponding advecting velocity is
\begin{equation}
\boldsymbol c_d
=
\nabla_d^\perp\mathfrak A_{i,d}
=
\boldsymbol c_d^{(0)}
+
\boldsymbol c_d^{(\mathrm{fr})},
\end{equation}
where
\begin{align}
\boldsymbol c_d^{(0)}
&=
\nabla_d^\perp
\left(
\frac{H_d}{\rho_0f_d}
\right),
\\
\boldsymbol c_d^{(\mathrm{fr})}
&=
-
\nabla_d^\perp
\left\{
\frac{
\sinh(2s_d)+\sin(2s_d)
}{
2\rho_0\alpha_d f_d
\left[\cosh(2s_d)+\cos(2s_d)\right]
}
\right\}.
\end{align}
Here
\begin{equation}
\nabla_d^\perp\psi
\equiv
\left(
\partial_{y_d}\psi,
-\partial_{x_d}\psi
\right).
\end{equation}

Because the frictional correction to the advecting streamfunction is a function of depth $H_d$ and Coriolis parameter $f_d$, it can only increase or decrease the flow along $f_d/H_d$ contours (not $H_d$ contours). 
For shallow depths, friction opposes the tendency to flow around $f_d/H_d$ contours: it retards the flow.

\end{document}  